\chapter{Neighborhood Parents, Availability, and Connection Descent}
\label{ch:neighborhood-parents}

The pointwise construction extends over a neighborhood when the local laws,
coarse channels, and geometric presentations satisfy explicit compatibility
conditions. The following construction supplies those conditions through a
common ambient law. It also distinguishes a parent requiring a fixed set of
children from a parent whose available children vary with context. These are
declared extensions of the pointwise datum, with different domains.
\status{DEFINITION}

\section{A cover without a common intersection}
\label{sec:rg-cover-domains}

For this example take \(\mathcal C=U=(-3/2,3/2)^2\), with
\(c=(x,y)\), and four agents with labels \(0,1,2,3\). This local indexing
replaces the generic population indexing only in this example. Define
\begin{equation}
\begin{aligned}
 I_0&=(-8/5,-1/10),& I_1&=(-8/5,4/5),\\
 I_2&=(-4/5,8/5),& I_3&=(1/10,8/5),\\
 U_i&=U\cap(I_i\times\R).
\end{aligned}
\label{eq:rg-cover-open-domains}
\end{equation}
These open sets and the contextual base are declared data; the construction
places no probability measure on the base. \status{DEFINITION}

\theoremheading{The overlap pattern}{thm:rg-cover-nerve}
The sets \(U_i\) cover \(U\). Their nonempty pairwise intersections have
indices \(01,02,12,13,23\), their nonempty triple intersections have indices
\(012,123\), and their fourfold intersection is empty. In particular,
\begin{equation}
 U_{012}=U\cap\{-4/5<x<-1/10\},\qquad
 U_{123}=U\cap\{1/10<x<4/5\}.
\label{eq:rg-cover-triples}
\end{equation}
This is an established cover with genuine triple compatibility conditions.
\status{ESTABLISHED}

\paragraph{Proof.}
An intersection of the intervals is nonempty exactly when its largest lower
endpoint is smaller than its smallest upper endpoint. This gives all the
listed intersections. The intervals \(I_1,I_2\) alone cover the first
coordinate of \(U\), proving the cover assertion. Finally
\(U_0\cap U_3=\varnothing\), which also excludes the fourfold intersection.
\(\square\)

Let \(a_i,r_i\) be the center and half-width of \(I_i\). The availability
weights are the prescribed functions
\begin{equation}
 w_i(c)=
 \begin{cases}
 \exp(1-1/t_i(c)),&t_i(c)>0,\\
 0,&t_i(c)\leq0,
 \end{cases}
 \qquad t_i(c)=1-\left(\frac{x-a_i}{r_i}\right)^2.
\label{eq:rg-cover-bumps}
\end{equation}
These weights encode availability, not posterior source probabilities.
\status{DEFINITION}

\theoremheading{Smooth availability with an explicit null event}{thm:rg-cover-availability}
Each \(w_i\) is smooth, lies in \([0,1]\), and is positive exactly on
\(U_i\). Add a null endpoint \(\bot\) and define the endpoint channel
\begin{equation}
 T_c(i,j)=w_i(c)\mathbf1_{j=i}+(1-w_i(c))\mathbf1_{j=\bot},
 \qquad T_c(\bot,j)=\mathbf1_{j=\bot}.
\label{eq:rg-cover-availability-channel}
\end{equation}
For any fixed finite number of directed records, conditionally independent
applications of this channel to their endpoints give a smooth normalized
channel on the complete record state. Consequently smooth ambient event
probabilities give smooth visible joint event probabilities, including at
agent entry and exit. This is an established property of the declared
availability mechanism. \status{ESTABLISHED}

\paragraph{Proof.}
Within its positive domain, every derivative of \(\exp(-1/t)\) is that
exponential times a polynomial in \(1/t\); all such products tend to zero
as \(t\downarrow0\). Extension by zero is therefore smooth. Since
\(0<t_i\leq1\) inside the interval, the stated range follows. Each row of
\(T_c\) sums to one. Products of the endpoint rows remain normalized,
and their entries and finite sums against smooth ambient probabilities are
smooth. Applying the channel to a joint law preserves any dependence not
removed by the channel; it does not replace correlated records by a product
of their marginals. \(\square\)

For each channel \(x\in\{b,m\}\), let \(\eta^x_{ij}(c)\) be the
marginal law of one visible directed record, including \(\bot\), and set
\begin{equation}
 \alpha_i^x(c)=\sum_j\eta^x_{ij}(c),\qquad
 b^x_{ij}(c)=\frac{\eta^x_{ij}(c)}{\alpha_i^x(c)}
 \quad\text{only when }\alpha_i^x(c)>0.
\label{eq:rg-cover-conditional-attention}
\end{equation}
Here \(b^b\) and \(b^m\) occupy the respective recognition attention
slots \(\beta\) and \(\gamma\); a generative posterior row is obtained
from its own event law by the same operation. The joint event law, occupancy,
and conditional row remain different objects. \status{DEFINITION}

If every real ambient receiver has positive mass, then
\(\alpha_i^x>0\) exactly on \(U_i\), because its visible occupancy equals
its ambient receiver mass times \(w_i\). On this domain the conditional
row is smooth. At zero occupancy the event law imposes no conditional row:
all its entries in that receiver row are zero, regardless of any proposed
conditional completion. Thus smoothness of \(\eta^x,\alpha^x\) supplies
no identifiable attention at zero occupancy. No null outcome is removed or
renormalized away in this established calculation.
\status{ESTABLISHED}

\section{Compatible laws and the two parent domains}
\label{sec:rg-cover-law-descent}

For the constructive specialization, adopt one smooth normalized ambient
joint \(\mathbb P_c(dY,do,de\mid X)\), a smooth selected recognition
\(\mathbb Q_{c,o,X}(dY,de)\), and a reference law on the same finite
product of agent fibers and record states. The fixed structural datum is
\(X\). Assume positive finite evidence at the admitted observation and
selected posterior versions that are pushed with these same laws. The
continuous variables may be correlated Gaussian coordinates with positive
definite ambient covariance, while records have finite laws. An individual
agent's section is exposed only on its \(U_i\). These ambient laws are
additional coupling data, not a conclusion inferred from independently
specified local marginals. \status{HYPOTHESIS}

For agent subsets \(S\subset T\), let \(r_{S\leftarrow T}\) project
continuous coordinates to \(S\) and send every record endpoint outside
\(S\) to \(\bot\). Null endpoints remain null. Restriction of the base
variable to a smaller open set is a separate operation from this projection.
\status{DEFINITION}

\theoremheading{Compatibility supplied by an ambient extension}{thm:rg-cover-compatible-laws}
Every ambient probability law in the preceding hypothesis induces compatible
local laws. On nested subsets and open context domains,
\begin{equation}
 r_{R\leftarrow S}r_{S\leftarrow T}=r_{R\leftarrow T},
 \qquad
 (\mathbb Q_T r_{S\leftarrow T})r_{R\leftarrow S}
 =\mathbb Q_T r_{R\leftarrow T}.
\label{eq:rg-cover-restriction-law}
\end{equation}
The same equality holds for the generative joint, reference, and selected
posterior families. Smooth local presentations that obey their transition
maps on pairwise overlaps obey the corresponding triple-overlap cocycles.
This is an established sufficient construction of compatibility.
\status{ESTABLISHED}

\paragraph{Proof.}
Coordinate projections compose. Sending a removed endpoint to \(\bot\)
and retaining \(\bot\) also composes. Pushforward through these deterministic
channels therefore gives the displayed identity. Restricting the same
formula in \(c\) twice is its restriction to the smaller domain. If
\(G_a,G_b,G_d\) are three local frame matrices for one associated fiber,
the coordinate transformations are \(G_bG_a^{-1}\), and
\((G_dG_b^{-1})(G_bG_a^{-1})=G_dG_a^{-1}\). Pushing all coordinates and
full covariances by those same transformations proves the law cocycle.
These are presentations of one law, not equality of different agents'
beliefs. \(\square\)

A general extension theorem for arbitrary local laws is not asserted here.
Its unresolved obligation is to characterize which full local probability
objects, supplied without the ambient extension, admit one joint extension
compatible with every restriction and selected conditional version. The
present proof starts with that extension and therefore cannot establish the
general implication. \status{OPEN}

For a fixed block \(I_A\), the fixed-child domain and the availability domain
are, respectively,
\begin{equation}
 U_A^{\cap}=\bigcap_{i\in I_A}U_i,\qquad
 V_A^{\cup}=\bigcup_{i\in I_A}U_i,
 \qquad s_A(c)=\sum_{i\in I_A}w_i(c).
\label{eq:rg-cover-parent-domains}
\end{equation}
The original \(U_A\) of the pointwise construction means
\(U_A^{\cap}\). On \(V_A^{\cup}\), declare the different availability
parent by coefficients \(\lambda_{Ai}=w_i/s_A\), using all the locally
available joint child data supplied by the ambient extension.
\status{DEFINITION}

For the explicit geometric example use the standard two-dimensional
representation of \(\mathrm{SO}(2)\) separately for belief and model, and
let \(G_{\mathrm f},G_A\) be the resulting orthogonal frame matrices for
the active fine and parent coordinates. Write
\begin{equation}
 L_A=(\lambda_{Ai}I_4)_{i\in I_A},\qquad
 F_A=G_A L_A G_{\mathrm f}^{-1}.
\label{eq:rg-cover-parent-channel}
\end{equation}
This is a deterministic normalized channel on the continuous coordinates.
Its event channel sends a child endpoint to its block label and keeps
\(\bot\). This declared example assumes equal standard representations;
an inequivalent pair of representations requires its own intertwiners.
\status{HYPOTHESIS}

\theoremheading{Availability parents glue and compose}{thm:rg-cover-parent-gluing}
The declared availability parent is smooth on \(V_A^{\cup}\) and agrees
on all chart overlaps. For disjoint blocks, the active rows of their combined
\(F\) are independent, so a positive definite fine Gaussian covariance
pushes to a positive definite covariance on the active parent fiber. If a
second stage weights the active blocks by \(s_A/\sum_Bs_B\), its direct
fine coefficient is \(w_i/\sum_jw_j\). Thus the direct and staged complete
law channels agree. This is an established extension with changing child
availability, rather than continuation of a fixed-child object by necessity.
\status{ESTABLISHED}

\paragraph{Proof.}
On \(V_A^{\cup}\), \(s_A>0\), so every coefficient is smooth. At a
child boundary its vanishing bump extends smoothly as long as another child
keeps the denominator positive. The formula is the same in every context
chart, with the frame cocycle already proved. Each active block row has
nonzero coefficients on its own disjoint child block, which gives full row
rank. For a nonzero parent vector \(v\),
\(v^\top F\Sigma F^\top v=(F^\top v)^\top\Sigma(F^\top v)>0\).
At two stages the coefficient product is
\begin{equation}
 \frac{s_A}{\sum_Bs_B}\frac{w_i}{s_A}
 =\frac{w_i}{\sum_jw_j}\quad(i\in I_A, s_A>0).
\label{eq:rg-cover-staged-weights}
\end{equation}
An absent block contributes zero directly; no conditional coefficient for
that absent block is needed. Frame matrices cancel at intermediate fibers.
Endpoint maps compose with the same null convention, and equality of complete
channels gives equality of all their pushed laws. For blocks \(01\) and
\(23\), the final availability parent exists on all \(U\), whereas a
fixed four-child parent has empty domain by \cref{thm:rg-cover-nerve}.
\(\square\)

The full fine Gaussian family is smooth on each constant-active-set region.
Across changing dimension, the construction does not supply one fixed-rank
Gaussian statistical manifold. Nor does it extend a parent through a point
where all its children are unavailable. These additional smooth-extension
claims are declined because their rank and domain premises fail in the
displayed construction. \status{NOT-CLAIMED}

\theoremheading{Evidence and the discarded conditional law}{thm:rg-cover-evidence}
At each fixed admitted context let \(C\) be any of the normalized,
recognition-independent state channels above, keeping the observation
unchanged. It preserves evidence and pushes the selected posterior version
to a posterior version of the pushed joint. For finite fine relative entropy,
\begin{equation}
 \mathcal F_{\mathrm f}=\mathcal F_{\mathrm c}
 +\int \KL\bigl(\mathbb Q(dY\mid z)\Vert
                  \boldsymbol\Pi(dY\mid z,o,X)\bigr)
          \mathbb Q_{\mathrm c}(dz).
\label{eq:rg-cover-vfe-disintegration}
\end{equation}
For a stochastic channel, \(Y\) denotes the original state in its joint
input-output extension. This established identity compares coarse levels at
fixed context and observation. \status{ESTABLISHED}

\paragraph{Proof.}
Integrating \(C(Y,dz)\) over \(z\) gives one, leaving the original
observation marginal. Integrating the selected original posterior against
\(C\) gives a conditional version of this pushed joint by its defining
integral identity. The common-channel input-output laws have the original
relative entropy because their conditional output kernel is identical.
Factor their density ratio into the retained marginal ratio and the
conditional input ratio, and integrate. This gives the KL chain rule;
adding the same finite \(-\log p(o\mid X,c)\) proves the displayed VFE
identity. Exceptional observation versions are transported by the chosen
version family, not inferred from an almost-sure identity at an arbitrary
observation. \(\square\)

The channel depends on context, so \cref{thm:rg-cover-evidence} does not
imply Fisher contraction for derivatives in \(c\). Such derivatives also
differentiate the channel. Nor does a comparison across coarse levels give
an inference-time decrease or increased evidence. These implications are
declined because they compare different derivatives or different objects.
\status{NOT-CLAIMED}

\section{Nonflat connections and metric-dependent projections}
\label{sec:rg-cover-connections}

On the trivial common principal \(\mathrm{SO}(2)\) bundle over \(U\),
choose two connections independently, with local reference potentials
\begin{equation}
 A^x_*={\kappa_x\over2}(x dy-y dx)J,\qquad
 J=\begin{pmatrix}0&-1\\1&0\end{pmatrix},\qquad x\in\{b,m\}.
\label{eq:rg-cover-nonflat-potential}
\end{equation}
In this equation only, the superscript \(x\) labels a channel; the scalar
\(x\) inside the one-form is the first base coordinate. The convention is
\(\nabla=d+A\), so the potential in a frame \(G\) is
\(A=GA_*G^{-1}-dG G^{-1}\). The connections are supplied geometric data,
not quantities determined by attention. \status{DEFINITION}

\theoremheading{A nonflat parent retains its connection defect}{thm:rg-cover-nonflat-defect}
The reference curvatures are
\(\mathcal R^x_*=\kappa_xJ dx\wedge dy\). On a constant-rank parent
domain, with the same channel connection repeated in its child fibers,
\begin{equation}
 \mathcal D F=dF+A_{\mathrm c}F-FA_{\mathrm f}
   =G_{\mathrm c}(dL)G_{\mathrm f}^{-1}.
\label{eq:rg-cover-connection-defect}
\end{equation}
For composable maps their defects satisfy
\begin{equation}
 \mathcal D(F_2F_1)=(\mathcal D F_2)F_1+F_2(\mathcal D F_1).
\label{eq:rg-cover-defect-composition}
\end{equation}
These established equations permit nonzero curvature and do not make
variable-weight coarse maps parallel. \status{ESTABLISHED}

\paragraph{Proof.}
The abelian potential has \(A_*\wedge A_*=0\), and differentiating its
coefficient gives \(\kappa_xJ dx\wedge dy\). Differentiate
\(F=G_{\mathrm c}LG_{\mathrm f}^{-1}\) and substitute the displayed
frame potentials. Frame derivatives cancel. The repeated reference
connection commutes through scalar block weights, leaving \(dL\).
The product rule gives the second equation by adding and subtracting
\(F_2A_1F_1\). Thus gauge covariance is compatible with a nonzero
first-order defect. \(\square\)

\theoremheading{An operational record sensitive to a declared base holonomy}{thm:rg-cover-operational-holonomy}
Fix a probe agent at \(c=(7/20,1/5)\), and the closed rectangular path
with successive vertices
\((7/20,1/5)\), \((13/20,1/5)\), \((13/20,1/2)\), \((7/20,1/2)\), and \(c\).
Choose an agent whose domain contains this rectangle, a selector \(P_{\rm pr}\)
onto its separately typed belief and model coordinates, and fixed positive
definite noise \(R_{\rm pr}\). Let a finite label \(a\) specify prescribed
connection parameters \((\kappa_{b,a},\kappa_{m,a})\), with positive prior
masses \(\pi_a\). Keep one original normalized joint
\(p_0(Y,e,o_0)\), independent of \(a\). The additional record and joint
are declared by
\begin{equation}
\begin{aligned}
 U_a&=G(c)\exp\left[-\mathcal A 
       \operatorname{diag}(\kappa_{b,a}J,\kappa_{m,a}J)\right]G(c)^{-1},
 \qquad \mathcal A=9/100,\\
 H_a&=(U_a-I_4)P_{\rm pr},\\
 p(a,Y,e,o_0,r)&=\pi_a p_0(Y,e,o_0)
                       N(r;H_aY,R_{\rm pr}).
\end{aligned}
\label{eq:rg-cover-loop-observation}
\end{equation}
This is a normalized statistical experiment whose record can distinguish
base holonomies while the matter law, loop, selector, and noise are held
fixed. It is an established operational witness for this declared tuple.
\status{ESTABLISHED}

\paragraph{Proof.}
The rectangle lies inside \(U_2\); its convexity and the strip's convexity
keep every segment in that agent domain. Integrating the reference
one-form along a segment gives
\(\kappa\det(c_j,c_{j+1})/2\). Summing segments gives
\(\kappa\mathcal A\), and the equation \(\dot v=-A(\dot\gamma)v\)
gives the displayed exponential. Endpoint frames conjugate it at the
basepoint. Each new Gaussian factor integrates to one in \(r\), then
\(p_0\) and \(\pi\) normalize the full joint.

If the original matter posterior is \(N(\mu_0,\Sigma_0)\), with
\(\Sigma_0\) positive definite, Gaussian integration gives
\begin{equation}
 \ell_a(r)=N(r;H_a\mu_0,R_{\rm pr}+H_a\Sigma_0H_a^\top),\qquad
 Z(r,o_0)=Z_0(o_0)\sum_a\pi_a\ell_a(r),\qquad
 \omega_a=\frac{\pi_a\ell_a(r)}{\sum_d\pi_d\ell_d(r)}.
\label{eq:rg-cover-loop-evidence}
\end{equation}
For example the supplied pairs \((0,0)\) and \((7/10,-2/5)\) give
\(H_0=0\) and \(H_1\ne0\), since neither nonzero angle is a multiple
of \(2\pi\). Therefore
\(H_1\Sigma_0H_1^\top\ne0\), so the two predictive Gaussian laws
have different covariance matrices. Their record laws differ with all the
listed non-holonomy data fixed. No particular observed value or unique
parameter recovery is required for this distinction.

Under a passive change of fine and probe coordinates by the represented
orthogonal frame maps \(Q,V\), send
\(P_{\rm pr}\) to \(VP_{\rm pr}Q^{-1}\), \(U_a\) to
\(VU_aV^{-1}\), \(r\) to \(Vr\), and \(R_{\rm pr}\) to
\(VR_{\rm pr}V^\top\). Then \(H_a'=VH_aQ^{-1}\), so the entire
observation kernel is pushed coherently and its evidence is invariant.
A zero-area loop instead gives \(H_a=0\) for every label and no label
reweighting; its record still has an ordinary common evidence factor.
\(\square\)

The witness declares both a connection and a curve; it does not identify an
arbitrary graph-link product with base transport. Distinct connection labels
can also have the same holonomy modulo a full rotation, and a chosen matter
law can leave further symmetries unidentifiable. A universal connection
recovery or graph-to-base identification is therefore not claimed.
\status{NOT-CLAIMED}

For a separate constant-rank construction, let \(E\) carry a declared
Euclidean fiber metric and a compatible connection \(d+A\), with
\(A^\top=-A\) in orthonormal frames. Let \(F:E\to E_{\mathrm c}\)
be a smooth surjective vector-bundle map of fixed rank. Define
\begin{equation}
 g=(FF^\top)^{-1},\qquad R=F^\top g,\qquad
 P=RF,\qquad N=I-P,
 \qquad B=F(dR+AR)=(FA-dF)R.
\label{eq:rg-cover-quotient-data}
\end{equation}
Here \(P\) is a fiberwise orthogonal projector, not a probability law;
\(g\) is the chosen metric's quotient metric. This is a declared
minimum-norm splitting and its projected connection.
\status{DEFINITION}

\theoremheading{Projected connection versus exact connection descent}{thm:rg-cover-projected-connection}
The connection \(d+B\) preserves \(g\). If
\(S=N(dR+AR)\) and \(\mathcal E=dF+BF-FA\), then
\begin{equation}
 dg=B^\top g+gB,\qquad
 \mathcal E=(dF-FA)N=g^{-1}S^\top.
\label{eq:rg-cover-projected-defect}
\end{equation}
Moreover \(F\) is parallel exactly when \(\ker F\) is preserved by the
fine connection. Its curvature satisfies
\begin{equation}
 \mathcal R_B=F\mathcal R_A R+g^{-1}S^\top\wedge S,
 \qquad \mathcal R_A=dA+A\wedge A.
\label{eq:rg-cover-projected-curvature}
\end{equation}
These identities hold for the declared metric and splitting of fixed rank. \status{ESTABLISHED}

\paragraph{Proof.}
Direct multiplication gives \(FR=I\), \(P=P^\top=P^2\), and
\(R^\top R=g\). A fine representative of \(y\) is uniquely
\(Ry+v\), with \(v\in\ker F\), and its squared norm is
\(y^\top gy+\|v\|^2\); this proves the quotient-metric description.
Multiplying \(B\) by \(g\) gives
\(gB=R^\top(dR+AR)\). Adding its transpose cancels the skew \(A\)
terms and gives \(dg\). Differentiating \(FR=I\) gives the formula
for \(\mathcal E\). With \(D_AP=dP+[A,P]\), differentiating
\(P^2=P\) shows that \(D_AP\) has only off-diagonal blocks, and
metric compatibility makes its coefficients symmetric. Hence
\((D_AP)R=S\) and
\(\mathcal E=F(D_AP)N=g^{-1}S^\top\). Its vanishing is equivalent
to preservation of the retained subbundle and, by differentiating orthogonal
inner products, of its orthogonal complement \(\ker F\).

On a retained section \(u=Pu\), the projected connection is \(P\nabla\).
Expanding its square yields
\(P\mathcal R_Au+P(D_AP)\wedge(D_AP)u\): the term with
\(P(D_AP)P\) vanishes, and \(N\nabla u=(D_AP)u\).
Transporting through \(R\) and using the preceding adjoint identities gives
\cref{eq:rg-cover-projected-curvature}. The second term is an
endomorphism-valued two-form, not a nonnegative scalar energy.
\(\square\)

The projected connection can be curved even with a flat source. For
\begin{equation}
 A=0,\qquad F(x,y)=\begin{pmatrix}1&0&-x\\0&1&-y\end{pmatrix},
 \qquad \mathcal R_B(0)=-J dx\wedge dy\ne0,
\label{eq:rg-cover-projected-curvature-example}
\end{equation}
this is an established two-dimensional quotient example.
\status{ESTABLISHED}

\paragraph{Proof.}
At the origin \(g=I_2\) and the only nonzero row of \(S\) is
\((-dx,-dy)\). Its exterior square is
\(\left(\begin{smallmatrix}0&1\\-1&0\end{smallmatrix}\right)
 dx\wedge dy\), which gives the formula by
\cref{thm:rg-cover-projected-connection}. All bundles in this example
are trivial, so the curvature does not establish nontrivial topology.
\(\square\)

For the scalar-weight parent \(L=(\lambda_i I_d)_i\), with a common
internal connection \(A_*\), the same construction instead gives
\begin{equation}
 \sigma=\sum_i\lambda_i^2,\qquad
 g=\sigma^{-1}I_d,\qquad
 B=A_*-\tfrac12d\log\sigma I_d,\qquad
 \mathcal R_B=\mathcal R_{A_*}.
\label{eq:rg-cover-scalar-weight-quotient}
\end{equation}
This is an established cancellation in this specialization.
\status{ESTABLISHED}

\paragraph{Proof.}
Insert \(R=\operatorname{col}_i(\lambda_i I_d)/\sigma\) in
\cref{eq:rg-cover-quotient-data}. The additional scalar one-form is exact,
its square is zero, and its mixed wedge terms with \(A_*\) cancel.
The quotient-orthonormal coordinates multiply \(y\) by
\(\sigma^{-1/2}\) and remove this scalar term from \(B\).
The original variable-weight defect need not vanish. \(\square\)

\theoremheading{Composition requires propagation of the quotient metric}{thm:rg-cover-quotient-composition}
For smooth surjections \(F:E_0\to E_1\) and \(G:E_1\to E_2\), start
with metric \(I\) on \(E_0\), give \(E_1\) its induced metric
\(g_1=(FF^\top)^{-1}\), and use that metric in the second projection.
Then the quotient metrics and projected connections agree for staged and
direct projection, whether or not either map is parallel. This is an
established composition theorem for the specified splittings.
\status{ESTABLISHED}

\paragraph{Proof.}
The second lift is
\(R_G=g_1^{-1}G^\top(Gg_1^{-1}G^\top)^{-1}\). Multiplication gives
\begin{equation}
 R_FR_G=(GF)^\top((GF)(GF)^\top)^{-1}=R_{GF}.
\label{eq:rg-cover-quotient-lift-composition}
\end{equation}
The final quotient metric has the same denominator. The product rule and
\(FR_F=I\) now give
\(G(d+B_1)R_G=GF(d+A_0)(R_FR_G)\), which is the direct connection.
Resetting the intermediate metric is a different construction: with
\(F=\left(\begin{smallmatrix}2&0&0\\0&1&0\end{smallmatrix}\right)\),
\(G=(1,x)\), and \(A_0=0\), reset metric \(I\) gives
\(-x dx/(1+x^2)\), whereas the direct connection is
\(-x dx/(4+x^2)\). This verifies why the metric must be retained.
\(\square\)

The quotient construction does not supply a canonical effective-support
quotient for arbitrary semidefinite edge weights, nor an intertwiner into an
arbitrarily prescribed associated representation. Those extensions require
constant-rank support data, compatible represented actions, and an admitted
quotient or bundle map. The displayed theorem closes the chosen metric
construction only. \status{OPEN}

\section{Actual base descent with inherited nonzero curvature}
\label{sec:rg-cover-base-descent}

Let \(B=U\times(-1,1)\), with extra coordinate \(z\), and let
\(f:B\to U\) be \(f(c,z)=c\). Choose a smooth passive frame
\(h(z)=\exp(\gamma zJ)\), and let \(H\) repeat this representation in
each separately typed associated coordinate. Declare the principal and
associated maps
\begin{equation}
 \Pi((c,z),g)=(c,h(z)^{-1}g),\qquad
 \Psi((c,z),v)=(c,H(z)^{-1}v).
\label{eq:rg-cover-base-maps}
\end{equation}
They cover the noninjective base map \(f\). This example removes an
explicitly redundant base coordinate; it does not average contexts.
\status{DEFINITION}

\theoremheading{Compatible section, law, and connection descent}{thm:rg-cover-base-descent}
Given smooth target sections, laws, and connections on their stated target
domains, define their source presentations by \(H\), and define
\begin{equation}
 A_{\mathrm s}=Hf^*A_{\mathrm t}H^{-1}-dH H^{-1},\qquad
 F_{\mathrm s}=H_{\mathrm c}f^*F_{\mathrm t}H_{\mathrm f}^{-1}.
\label{eq:rg-cover-base-pullback}
\end{equation}
Then sections, full laws, evaluators, and channels descend through
\(\Psi\), provided observations and structural data are also pulled back.
Both channel connections are compatible with descent, their curvature is
\(Hf^*\mathcal R_{\mathrm t}H^{-1}\), and the vertical coarse defect
is the conjugate pullback of the target defect. This is an established
nonidentity base descent that permits nonflat target connections.
\status{ESTABLISHED}

\paragraph{Proof.}
The principal map is equivariant since
\(h^{-1}(ga)=(h^{-1}g)a\). Substituting \(h^{-1}g\) in the target
principal connection form \(g^{-1}A_{\mathrm t}g+g^{-1}dg\) gives
\cref{eq:rg-cover-base-pullback}. Differentiating \(\Psi=H^{-1}\)
then gives \(d\Psi+f^*A_{\mathrm t}\Psi-\Psi A_{\mathrm s}=0\).
The curvature law follows by expanding \(dA_{\mathrm s}+A_{\mathrm s}
\wedge A_{\mathrm s}\); the derivatives of \(H\) cancel. Expanding
the covariant derivative of \(F_{\mathrm s}\) similarly leaves the
conjugate pullback of \(\mathcal D F_{\mathrm t}\).

For a target kernel \(C_c\), define the source kernel by undoing its input
frame, applying \(C_c\), and restoring its output frame. Evaluating a
target measurable set gives
\begin{equation}
 \Psi_{\mathrm c\#}(\mathbb Q_{\mathrm s}C_{\mathrm s})
 = (\Psi_{\mathrm f\#}\mathbb Q_{\mathrm s})C_{\mathrm t}.
\label{eq:rg-cover-base-channel-square}
\end{equation}
The same change of variables holds for the full generative joint, reference,
and selected posterior versions. Keeping observation coordinates unchanged
therefore preserves evidence; common invertible presentation maps preserve
relative entropy by cancellation of density Jacobians. The declared
normalized evaluator transforms by its separately typed input-output square.
Finally restrictions to \(U_i\) pull back to restrictions to
\(U_i\times(-1,1)\), so all cover identities persist. \(\square\)

For the reference potentials in \cref{eq:rg-cover-nonflat-potential},
\(A_{\mathrm s}^x=f^*A_*^x-\gamma J dz\); its mixed curvature is
zero and its inherited \(dx\wedge dy\) curvature can be nonzero. A
nonzero coefficient in the forgotten direction is therefore compatible with
exact descent. This is an established consequence of the prescribed passive
frame, not a new force in that direction. \status{ESTABLISHED}

\theoremheading{The descent criterion and its obstructions}{thm:rg-cover-descent-obstructions}
Let \(f:B\to\overline B\) be a given smooth surjective submersion and
\(\Psi:E\to\overline E\) a given smooth bundle map over \(f\).
A smooth section \(s\) has a unique descended section exactly when
\(\Psi s\) is constant on every complete fiber of \(f\). For a given
smooth surjective principal submersion \(\Pi\), a connection form descends
exactly when its pullbacks by the two projections of
\(P\times_{\overline P}P\) agree. These are established criteria for
given maps, not existence of a canonical quotient.
\status{ESTABLISHED}

\paragraph{Proof.}
Necessity follows by evaluating a descended object at two lifts of the same
base point or tangent vector. Conversely, a submersion has local smooth
sections: in local coordinates augment its full-rank map by complementary
source coordinates and apply the inverse function theorem. Composing
\(\Psi s\) with such a local section defines a smooth target section;
fiber constancy makes the definitions independent of the local section and
hence consistent on overlaps. Surjectivity gives uniqueness. For forms,
choose a local section \(j\) of \(\Pi\) and set
\(\overline\omega=j^*\omega\). The fiber-product equality applied to
\((v,dj d\Pi(v))\) proves
\(\omega=\Pi^*\overline\omega\), and also independence of \(j\).
Equivariance and reproduction of principal vertical generators descend by
surjectivity and the equivariant group-fiber isomorphisms.

The criteria have concrete failures. The source law
\(N(H(z)(1+z/4,0)^\top,I_2)\) descends under \(\Psi\) to a law
whose mean varies with \(z\); its KL to \(N(0,I_2)\) is
\((1+z/4)^2/2\), so no passive reframing makes the pair constant on the
forgotten fiber. Separately, adding \(\lambda zJ dx\),
\(\lambda\ne0\), to \(A_{\mathrm s}\) adds
\(\lambda J dz\wedge dx\) to its curvature. Its contraction with
\(\partial_z\) is nonzero, whereas every conjugate pullback curvature
annihilates that vector. This excludes compatible connection descent through
this \(f\). \(\square\)

Choosing a probability measure on a varying base fiber and averaging its laws
would define a new mixture operation. It requires a base prior or reference,
a joint base-fiber model, and the induced observation weighting. Neither
fiber constancy nor a canonical such measure follows from the cover. General
informative base aggregation, rank-changing quotient geometry, and a
probability law on smooth random sections remain open beyond the explicit
constructions above. \status{OPEN}
